% GENERATED FILE. Edit tamil/pronunciation-reference.md, then run npm run generate:books.

\chapter*{The Tamil script --- a reference}
\addcontentsline{toc}{chapter}{The Tamil script (reference)}
\markboth{Tamil script}{Tamil script}

A \textbf{reference}, not a chapter — a place to look things up. You are \emph{not} meant to read this before the lessons: the lessons teach reading through real words, introducing each letter as a word needs it. This page just gathers the system in one spot.

\section*{How the Tamil script works (the three facts the lessons teach)}

\begin{enumerate}
\raggedright
\setlength{\itemsep}{0.35em}
  \item \textbf{Consonants carry a built-in \textquotedblleft{}a\textquotedblright{}.} \ta{க} = \textquotedblleft{}ka,\textquotedblright{} \ta{ம} = \textquotedblleft{}ma,\textquotedblright{} \ta{ர} = \textquotedblleft{}ra.\textquotedblright{} You never write the short \emph{a} — it's already there.
  \item \textbf{A vowel sign changes that built-in vowel.} Hooked onto a consonant, it swaps the \emph{a} for another vowel — \ta{க} ka + \ta{ி} gives \ta{கி} ki. The signs are \textbf{\ta{ா}} = ā (\ta{கா} = kā) · \textbf{\ta{ி}} = i (\ta{கி} = ki) · \textbf{\ta{ீ}} = ī · \textbf{\ta{ு}} = u · \textbf{\ta{ூ}} = ū · \textbf{\ta{ெ}} = e · \textbf{\ta{ே}} = ē · \textbf{\ta{ை}} = ai · \textbf{\ta{ொ}} = o · \textbf{\ta{ோ}} = ō · \textbf{\ta{ௌ}} = au.
  \item \textbf{A \emph{puḷḷi} (the dot, \ta{்}) removes the built-in vowel.} \ta{க்} = bare \textquotedblleft{}k.\textquotedblright{} A vowel-less consonant then leans into the next (\ta{க்} + \ta{க} $\to$ \ta{க்க}, a held \textquotedblleft{}kk\textquotedblright{}).
\end{enumerate}

\section*{Independent vowels (word-initial)}

When a word \emph{starts} with a vowel, a full vowel letter is used: \textbf{\ta{அ}} a · \textbf{\ta{ஆ}} ā · \textbf{\ta{இ}} i · \textbf{\ta{ஈ}} ī · \textbf{\ta{உ}} u · \textbf{\ta{ஊ}} ū · \textbf{\ta{எ}} e · \textbf{\ta{ஏ}} ē · \textbf{\ta{ஐ}} ai · \textbf{\ta{ஒ}} o · \textbf{\ta{ஓ}} ō · \textbf{\ta{ஔ}} au.

\section*{What makes Tamil distinctive}

\begin{itemize}
\raggedright
\setlength{\itemsep}{0.35em}
  \item \textbf{Three \emph{n}'s, three \emph{l}'s, two \emph{r}'s.} Tamil \emph{writes} the place of articulation: dental \textbf{\ta{ந}} / alveolar \textbf{\ta{ன}} / retroflex \textbf{\ta{ண}}; \textbf{\ta{ல} \ta{ள} \ta{ழ}}; \textbf{\ta{ர} \ta{ற}}. English fuses each set into one letter and never notices; in Tamil they are distinct and can change meaning.
  \item \textbf{One letter, several sounds.} Tamil marks no voicing or aspiration: \textbf{\ta{க}} can be \emph{k}, \emph{g}, or \emph{h}; \textbf{\ta{ச}} can be \emph{s} or \emph{ch} — \textbf{position} decides. The alphabet is small; the exact sound is read from context, picked up by ear.
  \item \textbf{No Sanskrit-only letters in the core set.} Classical Tamil gets by with 18 consonants; a handful of \emph{grantha} letters (\textbf{\ta{ஜ}} ja, \textbf{\ta{ஷ}} ṣa, \textbf{\ta{ஸ}} sa, \textbf{\ta{ஹ}} ha) are added only to spell loanwords.
\end{itemize}

\section*{Tamil in the Dravidian family}

\textbf{Dravidian} is a language family entirely separate from Indo-European (English, Latin, Sanskrit). Tamil is its oldest-attested member and the most protective of its native word-stock; its sisters — \textbf{Malayalam} (closest), then \textbf{Kannada} and \textbf{Telugu} — share the grammar and many roots but borrowed far more freely from Sanskrit. The lessons trace, word by word, which Tamil words are home-grown and which the sisters replaced (\emph{vaṇakkam} vs. \emph{namaskāra}; \emph{naṉṟi} vs. \emph{dhanyavāda}).
